\section{Manual de Referencia Técnica MDA-Audio-PIM-Relations-Machines}

Este lenguaje, \textbf{MDA-Audio-PIM-Relations-Machines}, permite orquestar las interconexiones entre máquinas PIM. A diferencia del nivel CIM, donde las relaciones son entre máquinas completas, el nivel PIM establece vínculos directos entre puntos de conexión y parámetros específicos.

\section{Estructura del Documento}

El lenguaje utiliza una estructura JSON para definir el grafo de relaciones.

\subsection{Propiedades Raíz}

\begin{table}[H]
\centering
\small
\begin{tabular}{@{}llp{5cm}l@{}}
\toprule
Atributo & Tipo & Descripción & Restricciones \\ \midrule
\texttt{description} & STRING & Propósito de la red de relaciones. & Opcional. 1 a 600 car. \\
\texttt{relations} & ARRAY & Lista de objetos de relación detallada. & Obligatorio. Puede estar vacío. \\ \bottomrule
\end{tabular}
\end{table}

\textbf{Ejemplo de Estructura Raíz:}
\begin{lstlisting}[language=json]
{
  "description": "0123456789 (hasta 600 caracteres)",
  "relations": [
    {
      "id": "550e8400-e29b-41d4-a716-446655661001",
      "source": "550e8400-e29b-41d4-a716-446655445001",
      "destination": "550e8400-e29b-41d4-a716-446655444004",
      "description": "Modulacion de amplitud (10 a 300 caracteres)"
    }
  ]
}
\end{lstlisting}

\subsection{Objeto Relation (Relación PIM)}

Define un vínculo técnico entre un punto de salida y un punto de entrada o parámetro.

\begin{table}[H]
\centering
\small
\begin{tabular}{@{}llp{4cm}l@{}}
\toprule
Atributo & Tipo & Descripción & Restricciones \\ \midrule
\texttt{id} & STRING & Identificador único de la relación. & Obligatorio. UUID de 36 car. \\
\texttt{source} & STRING & ID del \textbf{output} de una máquina PIM. & Obligatorio. Debe ser externo. \\
\texttt{destination} & STRING & ID del \textbf{input} o \textbf{parámetro}. & Obligatorio. Debe ser externo. \\
\texttt{description} & STRING & Descripción de la modulación o flujo. & Obligatorio. 10 a 300 car. \\ \bottomrule
\end{tabular}
\end{table}

\textbf{Ejemplo de Relación Detallada:}
\begin{lstlisting}[language=json]
{
  "id": "550e8400-e29b-41d4-a716-446655661001",
  "source": "550e8400-e29b-41d4-a716-446655445001",
  "destination": "550e8400-e29b-41d4-a716-446655444004",
  "description": "Conexion de oscilador a ganancia del amplificador."
}
\end{lstlisting}

\section{Reglas de Validación}

El compilador de PIM Relations aplica las siguientes reglas de integridad:

\begin{enumerate}
    \item \textbf{Formato UUID Estricto:} Los campos \texttt{id}, \texttt{source} y \texttt{destination} deben tener exactamente 36 caracteres. No se permiten nombres descriptivos; solo identificadores técnicos.
    \item \textbf{Unicidad de IDs:} El \texttt{id} de cada relación debe ser único en todo el documento.
    \item \textbf{Restricciones de Texto:}
    \begin{itemize}
        \item Descripción del documento: máx. 600 caracteres.
        \item Descripción de relación: entre 10 y 300 caracteres.
    \end{itemize}
\end{enumerate}

\section{Guía de Uso Táctico}

Este lenguaje es el puente final antes de la generación de código de bajo nivel. Mientras que en \textit{CIM Relations} decimos \textit{``La Máquina A se conecta a la B''}, en \textit{PIM Relations} especificamos \textit{``El output\_1 de la instancia 550e8400... se conecta al frequency de la instancia a1b2c3d4...''}.

\begin{tcolorbox}[colback=red!5!white,colframe=red!75!black,title=IMPORTANTE]
Los IDs utilizados en \texttt{source} y \texttt{destination} deben corresponder a identificadores definidos en los archivos \texttt{.json} de máquinas PIM correspondientes, siguiendo las reglas ya expuestas.
\end{tcolorbox}
